\clearpage
\hypertarget{native-479-lower}{}
\begin{center}
\rule{0.55\linewidth}{0.4pt}\\[2pt]
Astronomische Nachrichten. Band XXVI. Nr. 613. S. 193..208. 14. Oktober 1847.\\[2pt]
\rule{0.55\linewidth}{0.4pt}
\end{center}
\begin{center}[Beobachtungen der \emph{Iris} am Meridiankreis von \textsc{Gauss}.]\end{center}
\begin{flushright}Göttingen, 5. Sept. 1847.\end{flushright}

Herr \textsc{Gould} hat mir soeben die Elemente geschickt, die er aus meinen Beobachtungen der Iris vom 21. und 28. August zusammen mit der Beobachtung aus Cambridge vom 14. August abgeleitet hat. Ich beeile mich, sie Ihnen mitzuteilen. --- --- --- ---

Gestern Abend habe ich Iris noch einmal am Meridiankreis beobachtet.

\begin{center}
\setlength{\tabcolsep}{4pt}
\begin{tabular}{l|r|r|r}
 & Göttinger M. Z. & AR. & Deklination \\\hline
1847 Sept. 4 & $8^{\rm u}\ 51^{\rm m}\ 12^{\rm s}0$ & $296^\circ\ 7'\ 14''1$ & $-14^\circ\ 6'\ 3''8$ \\
\hfill 6 & $8\ 42\ 54.5$ & $296\ 0\ 46.5$ & $-14\ 8\ 56.1$ \\
\hfill 10 & $8\ 26\ 43.6$ & $295\ 53\ 59.8$ & $-14\ 13\ 59.5$ \\
\end{tabular}
\end{center}

Der Planet erschien bei dieser Beobachtung lichtschwächer als bei jeder früheren. Wenn der Planet zu sehen sein sollte, blieben die Fäden ganz undeutlich. Das mag zu einem großen Teil am Zustand der Luft gelegen haben, denn alle Sterne erschienen (mit bloßem Auge) in mattem, diesigem Licht, obwohl der Himmel ganz wolkenfrei war.

\begin{center}
\begin{tabular}{l|r|r|r}
Sept. 12 & $8^{\rm u}\ 18^{\rm m}\ 50^{\rm s}5$ & $295^\circ\ 53'\ 37''8$ & $-14^\circ\ 16'\ 16''3$ \\
\end{tabular}
\end{center}

Din Luft war an diesem Abend günstiger als am 10. Am 13. dagegen war Iris nur im ganz verdunkelten Gesichtsfeld gerade noch zu sehen. Eine Beobachtung war unmöglich. Der etwa $1^{\rm m}\ 40^{\rm s}$ vorausgehende Stern 7. 8. Größe, der auf fast demselben Parallelkreis stand, hatte sich noch bei dürftiger Beleuchtung beobachten lassen.

\begin{center}
\begin{tabular}{l|r|r|r}
Sept. 19 & $7^{\rm u}\ 52^{\rm m}\ 19^{\rm s}9$ & $296^\circ\ 8'\ 51''1$ & $-14^\circ\ 22'\ 30''0$ \\
\end{tabular}
\end{center}
